% ID: SEL-D0006-OB-E04
% ESTADO: cotejada. Segunda pasada visual completa y borrador compilado revisados.
% Fuente: Selectividad/sel_2019_matematicas_aplicadas/Reserva B/Exámenes Matemáticas Aplicadas CCSS - A y B.pdf
% SHA-256: f9236ec66af0dbd3fbd0231fdcaa6d922425c04d48178ea0737073e76e8eb7d5
% Página: 2; opción B; ejercicio 4.

En una muestra de 320 personas jubiladas elegidas al azar en un distrito de
una ciudad, resultó que 96 de ellas realizaban alguna actividad física.
\begin{enumerate}[label=\alph*)]
  \item \textbf{(1.5 puntos)} Construya un intervalo de confianza al
  95\,\% para la proporción de personas jubiladas que realizan alguna
  actividad física en ese distrito.
  \item \textbf{(1 punto)} Suponiendo que se mantiene la misma proporción
  muestral, halle el tamaño mínimo de la muestra para que el error cometido
  sea inferior a $0.1$ con un nivel de confianza del 98\,\%.
\end{enumerate}
